
\documentclass[11pt]{article}

\usepackage{amsmath,amssymb}
\usepackage{geometry}
\geometry{margin=1in}

\title{A Data-Driven Framework for Learning the Dynamics of Many-Body Microfluidic Systems}
\author{}
\date{}

\begin{document}
\maketitle

\tableofcontents
\newpage

\section{Introduction}

The collective dynamics of droplets flowing through microfluidic networks 
constitute a prototypical nonequilibrium many-body system. Although the 
governing equations of fluid motion are known, the effective dynamics 
emerging from hydrodynamic interactions, confinement, droplet deformation, 
and evolving pressure fields are highly nonlinear and history-dependent. 
Consequently, deriving closed-form evolution equations for the observable
state of interacting droplets remains a formidable challenge. Existing
descriptions therefore rely either on simplified physical approximations 
or on device-specific rules whose applicability is often restricted to 
particular geometries and operating conditions.
Instead of deriving the effective equations of motion analytically, 
we seek to infer the time-evolution operator directly from experiment. 
Ultimately, the objective is not merely to forecast droplet trajectories, 
but to construct an effective dynamical description of an interacting 
nonequilibrium system directly from experiment. In this sense, the 
learned model functions as a data-driven approximation of the system's 
propagator, providing a route toward discovering effective state variables 
and interaction laws that are difficult—or perhaps impossible—to derive 
analytically from first principles.

\section{Scientific Objective}

The objective of this project is \textbf{not trajectory prediction.}

Prediction accuracy serves only as experimental validation that the learned
latent dynamics capture the governing physics.

The true objective is to

\begin{itemize}
\item approximate the unknown evolution operator,
\item discover latent physical state variables,
\item identify mechanisms governing collective droplet behavior,
\item build a surrogate dynamical model capable of scientific interpretation.
\end{itemize}

\section{Physical Formulation}

The underlying hypothesis is that the evolution of the observable state 
is governed by an unknown propagator,
\[
\mathbf{x}_{t+1}=\mathcal{T}(\mathbf{x}_t),
\]
where $\mathbf{x}_t$ denotes the complete microscopic state of the 
droplet ensemble at time $t$, and $\mathcal{T}$ represents the effective 
dynamics resulting from the combined action of hydrodynamic interactions, 
confinement, and geometric constraints. Recovering this operator from 
experimental observations constitutes the central objective of the present work.
Since the full microscopic state of the system is not experimentally 
accessible, we instead observe only a subset of physically measurable 
variables. Each droplet is represented by its position, velocity, and 
morphological descriptors, which together define the observable state of the system. These experimentally measured quantities are organized into a canonical spacetime representation that preserves the identity and temporal evolution of every physical droplet throughout its lifetime. Importantly, this representation remains intentionally close to the raw observables and avoids transformations into reduced coordinates or geometry-dependent variables, thereby minimizing assumptions regarding the structure of the underlying dynamics.
The observable state alone is generally insufficient to render the 
dynamics Markovian. Two experimentally indistinguishable configurations 
may evolve differently because they differ in unresolved variables, 
including local pressure fields, flow perturbations generated by neighboring 
droplets, or hydrodynamic memory accumulated through previous interactions. 
Consequently, the effective state of the system must be regarded as
\[
\mathbf{z}_t=\left(\mathbf{x}_t,\mathbf{h}_t\right),
\]
where $\mathbf{x}_t$ denotes the experimentally observable degrees 
of freedom and $\mathbf{h}_t$ represents latent variables encoding 
the unresolved physics. Rather than prescribing these hidden variables 
through phenomenological models, they are inferred directly 
from experimental data as part of the learning process.
The evolution operator is approximated using a Transformer 
architecture that treats the droplet ensemble as an interacting 
many-body system. Through self-attention, the model learns how the 
future evolution of each droplet depends on the instantaneous configuration
of all other droplets, effectively approximating the interaction 
kernel governing the collective dynamics without requiring explicit 
specification of pairwise hydrodynamic forces or constitutive equations. 
In this formulation, predictive accuracy serves not as the ultimate objective, 
but as evidence that the learned latent state captures the essential physical 
information required to approximate the true propagator of the system.



\section{Observable State}

The current observable state is

\[
[x,y,v_x,v_y,c],
\]

where $c$ denotes circularity.

These variables are measurements of the physical state rather than the state
itself.
The present implementation characterizes droplet morphology through 
low-dimensional geometric descriptors, such as circularity, which provide 
a first approximation of droplet deformation. Future developments will 
replace or augment these handcrafted descriptors with learned visual 
embeddings extracted directly from droplet images. Such embeddings are 
expected to encode richer information regarding interfacial deformation, 
local stress distributions, and other morphological signatures that influence 
hydrodynamic interactions. As increasingly informative observables are 
incorporated into the state description, the estimated evolution operator 
is expected to converge toward a more faithful representation of the 
underlying physical dynamics.
One way to learn such morphological descriptors is to use 
convolutional neural networks (CNNs) to extract features from raw droplet images. 
The CNN can be trained in a self-supervised manner to learn representations that 
capture relevant physical characteristics, such as shape, size, and 
deformation patterns. These learned features can then be integrated into the 
observable state vector, enhancing the model's ability to predict future 
dynamics based on richer morphological information.

\[
\mathbf{z}_{\mathrm{shape}} = \mathrm{CNN}(\text{droplet image}),
\]
leading to the observable state
\[
[x, y, v_x, v_y, \mathbf{z}_{\mathrm{shape}}].
\]



\section{Canonical Dataset Philosophy}

The canonical dataset is the permanent digital representation of the
experiment. It is not merely a training dataset but the authoritative record
from which all downstream datasets are derived.

Its design principles include:

\begin{itemize}
\item permanent droplet identities,
\item interpolated missing detections,
\item recomputed velocities,
\item explicit existence masks,
\item complete reproducibility,
\item modular downstream processing.
\end{itemize}

(Note: explain the construction of the canonical dataset, including the mask
and its dimensions.)

\section{Neural Architecture}

The proposed learning architecture is designed to approximate the unknown evolution operator while respecting the physical structure of the problem. Rather than predicting trajectories directly, the network learns how the measured state of interacting droplets evolves through time.

The processing pipeline is

\[
\begin{aligned}
\text{Measured State}
&\rightarrow
\text{MLP Encoder}
\rightarrow
\text{Latent Embedding} \\
&\rightarrow
(\text{Time}+\text{Slot}+\text{Mask Embeddings}) \\
&\rightarrow
\text{Space--Time Transformer}
\rightarrow
\text{MLP Decoder}
\rightarrow
\text{Future Velocities}.
\end{aligned}
\]

\subsection*{Encoder}

An MLP encoder projects the measured variables into a higher-dimensional latent space. Since the inputs are already structured physical measurements rather than raw images, convolutional processing is unnecessary at this stage.

\subsection*{Time, Slot and Mask Embeddings}

Time embeddings encode temporal ordering, slot embeddings preserve the identity of each droplet throughout a training sample, and mask embeddings indicate whether a droplet physically exists at a particular time. Together these embeddings provide the Transformer with temporal, identity, and existence information.

\subsection*{Space--Time Transformer}

Self-attention allows the interaction neighborhood to be learned directly from data rather than prescribed analytically. This is particularly attractive for many-body droplet systems where the dominant interactions change continuously as droplets move through the network.


\end{document}
